% !TeX root = ../thuthesis-example.tex

% 中英文摘要和关键字

\begin{abstract}
本项目面向多方言学习需求，构建了一个集多模态知识处理、人工智能分析与个性化学习规划于一体的智能化学习平台，服务对象涵盖普通学习者与语言学研究者。平台的核心技术包括多模态方言音系知识图谱、多智能体协作框架、自适应多步检索生成模型以及音系自进化生成与校准模型。通过整合语音信号处理、自然语言理解与生成、跨模态信息检索等技术，平台能够对用户的录音进行精准的声学参数提取与标准发音匹配，并诊断发音差异。

在实际使用中，用户可通过多轮录音练习与系统互动，系统在云端快速完成声学分析与多维度评估，并基于方言知识图谱与专家规则生成结构化诊断结果及自然语言反馈。同时，多智能体框架中的不同智能体分别负责发音准确性分析、韵律与流畅度评估、学习路径优化等任务，确保反馈的多角度与高针对性。平台还可输出多模态学习资源，包括语谱图、声调曲线、口型视频等可视化内容，使学习者获得沉浸式的学习体验。

平台采用云端部署与跨平台前端架构，支持Web端与移动App的无缝访问，保障高并发条件下的稳定性与可扩展性。其创新性在于将方言学习的深度语言知识与现代人工智能技术深度结合，不仅能显著提升发音学习效果，也为中国方言的保护、传承与推广提供了可持续的技术支持。
  \thusetup{
    keywords = {方言音系知识图谱, 多智能体, 发音分析, 自适应学习, 自进化},
  }
\end{abstract}

\begin{abstract*}
This project targets the demand for multi-dialect learning , delivering an intelligent learning platform that integrates multi-modal knowledge processing, AI-driven analysis, and personalized learning path generation. The platform serves both general learners and linguistic researchers. Its core technologies include a multi-modal dialectal phonology knowledge graph, a multi-agent collaborative framework, an adaptive multi-step retrieval-generation model, and a phonology self-evolution and calibration model. By combining speech signal processing, natural language understanding and generation, and cross-modal information retrieval, the system can precisely extract acoustic parameters from user recordings, match them to standard pronunciations, and diagnose pronunciation deviations.

In practice, users engage in multiple rounds of recording exercises with real-time interaction. The cloud-based system performs rapid acoustic analysis and multi-dimensional evaluation, then generates structured diagnostic results and natural language feedback based on the dialect knowledge graph and expert-defined rules. Within the multi-agent framework, different agents are responsible for tasks such as pronunciation accuracy assessment, prosody and fluency evaluation, and learning path optimization, ensuring feedback that is multi-perspective and highly targeted. The platform also delivers multi-modal learning resources—such as spectrograms, tone curves, and articulatory videos—to provide learners with an immersive and comprehensive learning experience.

With cloud deployment and a cross-platform front-end, the system supports seamless access from both web and mobile applications, ensuring stability and scalability even under high concurrency. Its innovation lies in integrating in-depth dialectal knowledge with state-of-the-art AI technologies, significantly improving pronunciation learning outcomes while providing sustainable technical support for the preservation, transmission, and promotion of Chinese dialects.
  \thusetup{
    keywords* = {dialect phonology knowledge graph, multi-agent, Pronunciation Analysis, adaptive learning, self-evolution},
  }
\end{abstract*}
